\documentclass[tikz,border=8mm]{standalone}
\usepackage[american]{circuitikz}

\begin{document}
\begin{tikzpicture}[every node/.style={font=\sffamily\small}]
  % Master rails: VRAW y=9, GND y=0, TAIL y=2.
  \draw (0,9) -- (11,9) node[right] {VRAW};
  \draw (0,0) -- (11,0) node[right] {GND};
  \draw (11,0) node[ground] {};

  % Reference: R3 terminates at REF; DZ anode=GND and cathode=REF.
  \coordinate (REFnode) at (1.5,4);
  \draw (1.5,9) to[R,l=R3] (REFnode);
  \draw (1.5,0) to[zD,l_=DZ] (REFnode);
  \draw (REFnode) node[circ,fill=black,inner sep=1pt] {};
  \node[left] at (REFnode) {REF};

  % Differential NPN pair.  Bases are connected by their named networks.
  \node[npn] (Q1) at (5,4) {};
  \node[npn] (Q2) at (9,4) {};
  \node[below right=5pt] at (Q1) {Q1};
  \node[below right=5pt] at (Q2) {Q2};
  \draw (REFnode) -| (4.15,4) -- (Q1.B);
  \draw (11,4) node[right] {FB} -- (Q2.B);

  % Shared tail and current source.
  \draw (Q1.E) |- (5,2) -- (9,2) -| (Q2.E);
  \draw (7,2) to[isource,l_=I1] (7,0);
  \draw (7,2) node[circ,fill=black,inner sep=1pt] {};
  \node[below right=4pt] at (7,2) {TAIL};

  % PNP active load; Q3 is diode-connected and Q4 mirrors its base current.
  \node[pnp,xscale=-1] (Q3) at (5,7) {};
  \node[pnp] (Q4) at (9,7) {};
  \node[above=5pt] at (Q3) {Q3};
  \node[above=5pt] at (Q4) {Q4};
  \draw (Q3.E) |- (5,9);
  \draw (Q4.E) |- (9,9);
  \draw (Q1.C) -- (Q3.C);
  \draw (Q2.C) -- (Q4.C);

  % Diode-connected Q3 and Q4 mirror-base corridor.  Q3.C is the
  % collector junction; mirrorbase is the separate common-base junction.
  \coordinate (mirrorbase) at (7,6.15);
  \draw (Q3.C) -| (mirrorbase);
  \draw (Q3.B) -| (mirrorbase) -| (Q4.B);
  \node[circ,fill=black,inner sep=1pt] at (Q3.C) {};
  \node[circ,fill=black,inner sep=1pt] at (mirrorbase) {};
  \draw (Q2.C) -- ++(1.2,0) node[right] {N2};
  \draw (Q2.C) node[circ,fill=black,inner sep=1pt] {};
\end{tikzpicture}
\end{document}
